\section{Certified reflection and interval sweep direction}
\label{sec:orbit-direction}

The preceding coorientation test removes the orientation-double query on
some supplied surfaces. The remaining base query can still produce a large
trace. Its cost depends on a representation choice: the AHT implementation
always transmits through a rightmost carrier and removes a terminal suffix.
The same equivalence relation can be presented in the opposite order, so
this convention need not give an efficient trace on a particular input.

The interval API and the supplied-normal-surface API now accept an explicit
\code{sweep\_direction} option: \code{forward}, \code{reverse}, or
\code{wide}. The default remains \code{forward}, preserving every old
complete result and proof. The last option compares the initial terminal
carrier widths and reflects if the left one is wider. This is a scheduling
heuristic. It is useful on the layered meridian examples but loses on some
other genuine normal surfaces. The implementation therefore exposes it for
explicit use without promoting it to the default or claiming an adaptive
worst-case improvement.

\subsection{Reflection preserves the entire equivalence relation}

For a universe $\{0,\ldots,M-1\}$, let $R(x)=M-1-x$.
An inclusive pairing $[a,b]\leftrightarrow[c,d]$ becomes
\[
 [M-1-d,M-1-c]\ \leftrightarrow\ [M-1-b,M-1-a].
\]
This order keeps the stored domain to the left of the stored range. The
order-reversal flag is unchanged: conjugating an affine isometry by two
reflections preserves its slope, and swapping domain and range takes its
inverse, which also preserves that slope. Thus $R$ induces a bijection
between old and new orbits, including singleton orbits, static gaps and
components generated by overlapping reflections. For $M=0$ the map is
the unique empty bijection.

The producer preserves the caller's original ordered pairing list in the
certificate. If it reflects the working list, it records one initial
\code{reflect} event and then uses the ordinary AHT reductions. There is
no change to transmission, trimming, contraction, periodic merging or
truncation. Reflection does not increase the maximum bit length of the
universe size and endpoints, and leaves the classical polynomial orbit-counting bound
in the number of pairings and $\log(M+1)$ unchanged.

The independent checker applies the displayed endpoint formula to its own
current rows. It imports neither the direction selector nor the producer.
The event has exactly one field, its operation name; additional fields
are rejected. Since the formula is valid on any current universe, replay
also permits later reflections and pairs of cancelling reflections.
The maintained producer only needs the single initial event.

To preserve the interpretation of old proof formats, reflected traces use
version three for the Fine--Wilf merger threshold and version four for
the classical AHT threshold. Versions one and two retain their former
thresholds and reject reflection events. A new-version proof need not
actually use reflection; a legal old reduction remains legal under the
corresponding threshold. Unsupported-version mutation tests now use
version five. This changes no orbit mathematics in the old versions.

With $k$ interval pairings and a $B$-bit bound on the universe size and endpoints,
constructing and
checking a reflection takes $O(k)$ integer operations, $O(kB)$ elementary
bit work and $O(k)$ stored rows. The trace adds one constant-size event.
Neither producer nor checker visits the $M$ individual points. Cancellation
is polled while transforming rows. The operation is preparation, not an
AHT cycle; a local-event replay allowance nevertheless counts its event.
An incomplete computation returns no certificate or orbit count.

\subsection{A deliberately limited wider-end selector}

The selector scans the original pairings, ignores identities, and computes
the width remaining after the existing overlapping-reflection trim. Among
pairings incident to the smallest left endpoint it takes the widest;
among those incident to the largest right endpoint it does the same.
It reflects precisely when the left width is larger, keeping the historical
direction on ties or when there are no nonidentity pairings.

This describes the initial terminal carriers before periodic mergers.
Static-gap contraction does not change the width of an occupied interval.
However, the selector does not predict later mergers, transmitted rows,
or the suffix actually exposed by a carrier. A large initial carrier can
therefore be a poor predictor of the total search. The selector adds a
linear scan and makes a single initial choice; it does not switch midway
through a reduction and it has no guarantee relative to the better of the
two directions. That distinction matters for the proposed introspective
recognizer: a future switching policy needs a resource or progress bound
that this width comparison does not provide.

The normal-surface wrapper passes the chosen option independently to every
actual surface, double, boundary and optional boundary-cone query. Its
coorientation shortcut still uses the original geometric arc system.
Reflection inside an orbit proof changes neither those geometric parity
constraints nor the outer source fingerprint. Coordinate-divisor lifting
and boundary classification retain their existing checks and proof formats.
The current diagram recognizer's scheduling settings are unchanged.

\subsection{Compatibility, counterexamples and measurements}

The frozen baseline is revision \code{3310e1e84}. The research driver loads
its interval producer, interval checker, normal-surface wrapper and outer
checker separately, explicitly reconnecting the old wrapper to the old
interval producer and the old outer checker to the old interval checker.
Shared immutable pairing objects have the same validated representation.
This prevents a comparison from silently running the new search inside
the old wrapper.

The audit compares all 1,275 maintained normal vectors in both multiplicity
modes and both boundary-classification modes. All 5,100 default-forward
records equal the frozen implementation exactly, including proofs,
statistics and cycles. Each configuration also runs both explicit reverse
and wider-end choices, giving 10,200 comparisons with the same topology
and successful independent replay. The frozen component-level oracle
checks the expected counts. Another 1,000 random interval systems use
literal pointwise components as an independent oracle; both directions
under both merger rules give 4,000 matching comparisons, while 2,000
default records reproduce the old kernel exactly.

Focused tests also enumerate all two-pairing systems through four points,
exercise random larger systems, transport integers exceeding 20,000 bits,
check exact cycle and replay-event caps, test cancellation and malformed
proofs, and replay with producer search and direction selection disabled.
Two additional reflection events test the involution directly.
All 29 focused tests and all 1,255 maintained tests pass.
The layered meridian base query shows the effect of changing order. The
wider-end option selects the reversed direction on these inputs; event
counts include its extra reflection event.
\begin{center}
\small
\begin{tabular}{rrrrrrr}
\toprule
$N$ & forward events & wide events & fwd. trans. & wide trans. & fwd. bytes & wide bytes\\
\midrule
4 & 52 & 43 & 30 & 20 & 3511 & 2718\\
16 & 400 & 151 & 330 & 68 & 30914 & 10152\\
64 & 5392 & 583 & 5130 & 260 & 448949 & 49951\\
256 & 82960 & 2311 & 81930 & 1028 & 7072861 & 367462\\
\bottomrule
\end{tabular}
\end{center}


For $N=256$, transmissions fall from 81,930 to 1,028 and serialized base
proof bytes from 7,072,861 to 367,462. The base query still performs 259
cycles, compared with 261 before. Thus the large trace reduction comes
primarily from doing less work per cycle. The displayed samples do not
by themselves prove a linear family bound, and the implementation retains
full-list scans within each cycle. A smaller trace is not a proof of
linear running time.

The wider-end rule has a concrete contrary example in the corpus entry
\path{finite_figureEight_relabel_1:6}. On this supplied surface, it
reduces total cycles from 45 to 41 but increases total local events from
421 to 621 and proof bytes from 28,581 to 42,682. Both full certificates
are retained and independently replayed. This is an ordering counterexample,
not a failed topology computation or a knot verdict. It also shows why an
adaptive policy cannot infer lower proof-production work merely from
fewer AHT cycles.

On the 1,275-vector corpus with default multiplicity reduction and no
boundary classification, the wider-end rule shortens 160 event traces,
lengthens 359 and leaves 756 unchanged. Aggregate events increase from
182,925 to 189,742, while bytes increase from 12,861,213 to 13,315,621.
Always reversing is worse on this metric: 199,448 events and 13,921,942
bytes. These are exact deterministic counts, independent of timing noise.

The isolated timing comparison uses five randomized rounds, with a separate
same-implementation control for each arm. Every timed call includes full
normal-surface topology production, independent replay, proof serialization
and hashing; input construction is outside the timer. The old arm uses the
frozen forward producer and checker, while the new arm explicitly requests
the wider-end rule. Both arms retain the preceding coorientation shortcut.
All timed calls use default multiplicity reduction and omit the optional
boundary classification, which is tested separately in the audit.
All measured calls finish and their topology digests agree. Per-arm proof
digests, counts and sizes are stable across all measured and warm-up calls.
There are 240 measured calls and 48 warm-ups; counting the surfaces inside
the batch calls gives 25,720 measured and 5,144 warm-up surface analyses.
\begin{center}
\small
\begin{tabular}{lrrrrr}
\toprule
Supplied surface & forward ms & wide ms & paired factor & old A/A & new A/A\\
\midrule
Layered 1 & 0.939 & 0.939 & 0.984 & 0.993 & 1.001\\
Layered 16 & 11.166 & 8.474 & 1.317 & 1.011 & 0.997\\
Layered 64 & 114.741 & 50.027 & 2.304 & 0.999 & 0.994\\
Layered 256 & 1767.995 & 524.711 & 3.369 & 0.997 & 0.992\\
Layered 64, binary scale & 130.815 & 52.799 & 2.240 & 1.012 & 1.003\\
Empty & 0.645 & 0.688 & 0.939 & 0.934 & 1.041\\
M\"obius band & 1.241 & 1.248 & 0.999 & 1.004 & 1.018\\
Vertex-link disc & 2.589 & 2.654 & 0.985 & 0.987 & 1.016\\
Relabelled figure-eight surface & 17.711 & 20.591 & 0.851 & 0.985 & 1.004\\
Interior trefoil surface & 11.453 & 12.033 & 0.956 & 0.997 & 0.995\\
Genus-two coherent & 6.943 & 7.454 & 0.937 & 0.985 & 1.021\\
Entire 1,275-vector batch & 4999.519 & 5147.753 & 0.971 & 0.990 & 1.000\\
\bottomrule
\end{tabular}
\end{center}


The final row times all 1,275 surfaces in one batch. A paired factor greater
than one favors the wider-end rule. It is the median of paired ratios,
not the ratio of the two timing medians. The layered-256 gain is
$3.369\times$, while the figure-eight counterexample has factor $0.851$
(about 17 percent slower). The entire batch has factor $0.971$, about
3 percent slower, with control factors $0.990$ and $1.000$. Across all cases,
control medians range from $0.934$ to $1.041$; the widest variation is on
the tiny empty-surface case. These measurements include the
selector overhead and the larger proofs on its misses. They support making
the alternative available on suitable supplied surfaces, but do not support
changing the general default. The explicit \code{forward} choice continues
to reproduce the old behavior exactly.


Reproduce the audit and isolated paired measurements with
\path{../fast/normal_orbit_research/direction.py}, using \code{audit} or
\code{benchmark}, a required \code{--output} path and optional
\code{--rounds}. Full records, counterexample certificates, source hashes
and validation logs are retained under \path{data/orbit-direction-*}.
The experiment identifies a substantial avoidable trace cost and a
counterexample to a simple automatic remedy. It provides a checked alternate
execution order, while leaving the broader adaptive scheduling and normal
surface discovery problems unresolved. It does not prove a new general
subexponential or quasipolynomial unknot-recognition bound.
